The first result is descriptive. Exact-year emergence rates rise sharply with local scene
thickness. When a city has fewer than 10 active local bands, the exact-year emergence rate is
1.4 percent. When it has 40 or more, the rate rises to 4.0 percent. The same pattern appears for
multi-band depth: emergence rises from 1.4 percent with 0--1 multi-band musicians to 3.7 percent
with 10 or more. Figure \ref{fig:critical_mass} reports this critical-mass pattern directly.

\begin{figure}[htbp]
\centering
\includegraphics[width=0.92\textwidth]{scene_critical_mass_event_rates.png}
\caption{Exact-year emergence rates by local scene thickness and multi-band depth.}
\label{fig:critical_mass}
\end{figure}

The main regression result appears in the preferred exact-year fixed-effects transition
specification. The counts-only roll-5 design is the cleanest current main-table object because it
isolates the stable unconditional scene terms. The dependent variable is whether a city-genre cell
that has not yet emerged crosses the operational fifth-band threshold in the current year.
Table \ref{tab:preferred_results} reports the preferred paired presentation: a core counts-only
panel and a composition extension. The main paper read comes from Panel A. Because the lagged
predictors are z-scored, the coefficients are directly comparable on a common within-sample scale.
Local spawning flow enters at 0.0134, target-genre active bands at 0.0505, and target-genre
multi-band musicians at 0.0227. Relative to the sample event rate of 3.09 percent, those are
economically large associations on the linear-probability scale: a one-standard-deviation increase
in local spawning flow corresponds to about 1.3 percentage points higher transition probability, a
one-standard-deviation increase in target-genre active bands corresponds to about 5.1 percentage
points, and a one-standard-deviation increase in target-genre multi-band depth corresponds to
about 2.3 percentage points. The largest coefficient is therefore not generic city size, but
focal-niche thickness.

\begin{table}[htbp]
\centering
\caption{Preferred exact-year fixed-effects scene results}
\label{tab:preferred_results}
\small
\begin{tabular}{@{}lcc@{}}
\toprule
Predictor & Panel A: core counts-only & Panel B: composition extension \\
\midrule
Overall active bands & 0.0135*** & 0.0174*** \\
 & (0.0043) & (0.0044) \\
Nontrivial communities & 0.0062*** & 0.0008 \\
 & (0.0018) & (0.0020) \\
Overall multi-band musicians & -0.0223*** & -0.0244*** \\
 & (0.0028) & (0.0030) \\
Local spawning flow & 0.0134*** & 0.0224*** \\
 & (0.0017) & (0.0025) \\
Spawn share &  & -0.0066*** \\
 &  & (0.0012) \\
Target-genre active bands & 0.0505*** & 0.0488*** \\
 & (0.0015) & (0.0017) \\
Target-genre broker musicians &  & 0.0579*** \\
 &  & (0.0028) \\
Target-genre broker share &  & -0.0724*** \\
 &  & (0.0022) \\
Target-genre multi-band musicians & 0.0227*** & 0.0331*** \\
 & (0.0015) & (0.0027) \\
Target-genre multi-band share &  & -0.0195*** \\
 &  & (0.0023) \\
\midrule
$N$ & 171,653 & 171,653 \\
Events & 5,306 & 5,306 \\
Event rate & 0.0309 & 0.0309 \\
$R^2$ & 0.0422 & 0.0581 \\
Fixed effects & City-genre + year & City-genre + year \\
Clustering & City & City \\
\bottomrule
\end{tabular}

\vspace{0.4em}
\parbox{0.96\textwidth}{\footnotesize Notes: Standard errors in parentheses. All predictors are
one-year-lagged \texttt{roll-5} averages z-scored within the estimation sample, so coefficients are
comparable as within-sample percentage-point changes in emergence probability from a
one-standard-deviation increase in the predictor. *** $p<0.01$. Panel A is the active headline
specification. Panel B is retained as a background composition extension because the share terms
are substantively informative even though the main paper read comes from the core counts-only
panel. The negative coefficient on overall multi-band musicians should be read as a conditional
control result: holding focal-genre depth fixed, broader city-wide overlapping labor appears less
helpful than overlap concentrated within the target niche.}
\end{table}


The control pattern is also interpretable once the table is read conditionally rather than
literally. Overall active bands remain positive, but overall multi-band musicians enter negatively.
That sign is not a claim that worker overlap is bad for scenes. It is the opposite comparison.
Holding fixed focal-genre active bands and focal-genre multi-band depth, broader city-wide
multi-band overlap appears to proxy for labor spread across many niches rather than for depth in
the focal one. The paper's mechanism read is therefore about niche-specific recombination, not
generic scene-wide overlap. This is also why the positive target-genre multi-band coefficient is
the more substantively important labor-pool margin.

The direct threshold-response package is reported in Appendix Tables
\ref{tab:appendix_threshold_response} and \ref{tab:appendix_threshold_specificity}. The first
result is encouraging. The coefficient pattern is not pinned to the fifth-band cutoff: across
alternative emergence thresholds of three, four, six, and eight bands, local spawning flow,
target-genre active bands, and target-genre multi-band musicians all remain positive. The harder
risk-set exercises are more sobering. When I drop observations one band below the fifth-band
cutoff, all three headline terms remain positive but become much smaller. When I keep only
city-genre cells that are still well below the cutoff, the coefficients attenuate sharply and the
multi-band term is no longer distinguishable from zero. The disciplined interpretation is
therefore narrower than the most expansive version of the project. The paper is strongest on
late-stage local niche consolidation rather than on the earliest seed stage of scene formation.
A focal labor-pool specificity check still supports the recombination margin: conditional on
target-genre active bands and total target-genre active musicians, target-genre multi-band
musicians remain positive while total target-genre active musicians turn negative.

The measurement-validation package is reported in Appendix Tables
\ref{tab:appendix_measurement_validation} and \ref{tab:appendix_historical_sanity}. The current
preferred sample is already geography-clean on the project's audit rule: region-like and malformed
city labels are excluded upstream, so none survive into the live fixed-effects panel. Tightening
the sample further does not materially change the main coefficients. Dropping the small dense
labels flagged for closer inspection, restricting the sample to `1990` onward, and restricting to
the high-coverage countries all leave the three headline terms close to their baseline values.
The historical sanity cases are also reassuring. Operational emergence dates such as Birmingham
heavy metal in `1977`, Tampa death metal in `1987`, Bergen black metal in `1992`, and Gothenburg
melodic death metal in `1996` are at least broadly plausible. The Bay Area case is useful for a
different reason: its raw timing looks plausible too, but the label is explicitly excluded because
it names a wider unit rather than a city.

The coefficient pattern points to a specific view of local creative development. City-genre cells
are more likely to transition into operational scene status when three conditions hold. First,
cities generate new projects through local spinouts. Second, the focal genre already has a thicker
local stock of active bands. Third, musicians in that focal genre work across multiple local
projects, which is consistent with labor pooling and recombination within the niche.
Figure \ref{fig:coefficients} visualizes the current headline coefficients.

\begin{figure}[htbp]
\centering
\includegraphics[width=0.78\textwidth]{scene_preferred_mechanism_coefficients.png}
\caption{Preferred core exact-year fixed-effects coefficients on z-scored lagged predictors.}
\label{fig:coefficients}
\end{figure}

Background role-composition checks add little to this interpretation. Only guitarist share among
switchers shows any visible signal, and even that is weakly negative rather than positive
($-0.0013$, $p = 0.062$); drums, bass, vocals, and keyboards are all close to zero. I therefore
leave role composition out of the headline narrative. The stable margin is broader
overlapping-worker depth within the focal niche, not one privileged instrument class.

The remaining digital-era checks are secondary and do not change the paper's main read. A broad
period split suggests that local spawning flow and target-genre active-band stock attenuate after
`2005`, while the target-genre multi-band-musician coefficient remains stable. That is useful as a
descriptive nuance, but it is not central to the paper's contribution. I therefore leave the
broader internet-adoption and city-broadband probes to Appendix Section \ref{app:digital_checks}.
Those exercises are either too broad, too late, or too small to identify a local digital
mechanism cleanly, and none of them overturn the main scene result.

The contribution is reduced-form rather than causal. That is still enough to establish a clear
empirical pattern. Thicker local capability, stronger spinout flow, and deeper within-genre
overlapping experience are all associated with later scene emergence. Heavy metal is useful here
not because it is culturally exceptional, but because the setting makes local project overlap and
worker mobility unusually observable. In that sense, the paper speaks to a broader economic
question: how localized capability and recombination shape the emergence of new creative
varieties.
